When vision contradicts text, multimodal large language models (MLLMs) consistently favor text—even when images provide clear evidence otherwise. This bias poses risks for applications requiring visual grounding, yet its cause remains unclear. In this paper, we uncover a surprising finding: models often \emph{get it right initially}, forming correct vision-based predictions in their intermediate layers, before \emph{changing their minds} and favoring text in the final output. We call this ``\textbf{\textit{late-layer textual override}}". The visual information is encoded, it simply does not survive to the output. More intriguingly, we find that \emph{how} predictions change reveals \emph{whether} they're correct: 85\% of failures shift toward text, while 89\% of successes shift toward vision. This directional signature enables a simple but powerful intervention: when we detect a confident visual prediction being suppressed, we restore it. We propose CALRD (\textbf{\underline{C}}onflict-\textbf{\underline{A}}ware \textbf{\underline{L}}ayer \textbf{\underline{R}}eference \textbf{\underline{D}}ecoding), a training-free method that recovers overridden predictions at inference time. Experiments across five MLLMs of varying architectures demonstrate 4–9\% absolute improvements on conflict benchmarks while maintaining standard performance, without training or external knowledge. The method works by recovering what the model already knew but failed to preserve.
